% =============================================================================
% Capitolo 1 — Introduzione
% =============================================================================
% Prima stesura, da rifinire dopo aver definito la bib e lo stile con il relatore.
% Le \cite{} usano placeholder key coerenti con la bibliografia proposta:
%   - gama2014survey       -> Gama et al. 2014
%   - lu2019learning       -> Lu et al. 2019
%   - sethi2018reliable    -> Sethi & Kantardzic 2018
%   - baier2020detecting   -> Baier et al. 2020
%   - bifet2007adwin       -> Bifet & Gavaldà 2007 (SIAM)
%   - bifet2006adwin       -> Bifet & Gavaldà 2006 (technical report UPC)
%   - gof1994patterns      -> Gamma et al. 1994
%   - scipy2020            -> Virtanen et al. 2020 (SciPy)
%   - montiel2021river     -> Montiel et al. 2021 (river)
% =============================================================================

\chapter{Introduzione}
\label{ch:introduzione}

\section{Contesto e motivazione}
\label{sec:contesto-motivazione}

Il \textit{Machine Learning} è diventato negli ultimi due decenni una tecnologia
pervasiva nell'industria del software, con applicazioni che spaziano dalla
classificazione automatica di testi al riconoscimento di immagini, dalla
previsione della domanda commerciale al rilevamento di frodi finanziarie. La
disponibilità di dataset di grandi dimensioni, l'aumento della potenza di calcolo
e la diffusione di framework open-source hanno reso possibile la costruzione di
modelli predittivi sempre più accurati e complessi, aumentandone rapidamente il
grado di adozione in ambito produttivo.

Tuttavia, l'accuratezza raggiunta in fase di addestramento non garantisce
prestazioni analoghe una volta che il modello viene messo \textit{in produzione},
ovvero utilizzato per prendere decisioni operative su dati nuovi non presenti nel
dataset di addestramento. Uno dei problemi più insidiosi che affliggono i sistemi
di Machine Learning in produzione è il cosiddetto \textit{drift}: un
cambiamento nel tempo delle proprietà statistiche dei dati in ingresso, delle
predizioni prodotte dal modello, o della relazione tra ingresso e uscita.

Quando si verifica un drift, le prestazioni del modello tendono a degradare
progressivamente, spesso in modo silenzioso, senza che venga generato alcun
errore esplicito né alcuna segnalazione automatica da parte del sistema. La
conseguenza di un drift non rilevato può essere estremamente gravosa: modelli
utilizzati per approvare richieste di credito, per classificare comunicazioni
di clienti, per prevedere la manutenzione di impianti industriali o per
rilevare comportamenti anomali possono iniziare a produrre decisioni
sistematicamente errate, con impatti economici, reputazionali e in alcuni
contesti anche legali significativi. La letteratura scientifica ha documentato
numerosi casi di degradazione di modelli in ambiente produttivo e ha
evidenziato come la mancanza di un adeguato sistema di monitoraggio costituisca
uno dei principali fattori di rischio nell'adozione di sistemi di Machine
Learning in ambito industriale \cite{gama2014survey,lu2019learning}.

In risposta a queste criticità è emersa una disciplina denominata
\textit{MLOps} (\textit{Machine Learning Operations}), che si propone di
applicare i principi del DevOps ai sistemi di Machine Learning, includendo tra
i suoi obiettivi principali il monitoraggio continuo dei modelli in produzione,
la rilevazione automatica del drift e l'attivazione di procedure di
riaddestramento quando necessario. Costruire un sistema robusto di
\textit{drift monitoring}, capace di identificare tempestivamente cambiamenti
significativi nei dati e di guidare le azioni correttive conseguenti,
rappresenta oggi una delle sfide più rilevanti nell'ingegnerizzazione dei
sistemi di Machine Learning a livello industriale.


\section{Contesto aziendale}
\label{sec:contesto-aziendale}

Il presente lavoro di tesi è stato svolto in collaborazione con
\textit{Engineering Ingegneria Informatica S.p.A.}, una delle principali
aziende italiane operanti nel settore della trasformazione digitale, con
attività in numerosi ambiti strategici tra cui pubblica amministrazione,
energia, telecomunicazioni, sanità e finanza. All'interno dell'azienda, il
ricorso a modelli di Machine Learning per l'automazione di processi
decisionali e per l'analisi predittiva sta diventando sempre più diffuso,
con la conseguente necessità di garantire che questi sistemi mantengano nel
tempo prestazioni conformi alle aspettative sancite in fase di sviluppo.

La proposta di tesi, elaborata sotto la supervisione del tutor aziendale
Daniele Fakhoury, si inserisce in un più ampio programma di ricerca applicata
che l'azienda sta portando avanti sul tema dell'MLOps. Nello specifico, il
lavoro mira a fornire fondamenta metodologiche e implementative per un
servizio di drift monitoring riutilizzabile, che possa in prospettiva essere
integrato nei processi di gestione dei modelli in uso o in via di sviluppo.
Il carattere aziendale del progetto pone accento su tre aspetti: la robustezza
pratica delle soluzioni proposte, la loro estendibilità a scenari applicativi
diversi, e la compatibilità con le principali tecnologie MLOps oggi in uso
nell'industria del software.


\section{Obiettivi della tesi}
\label{sec:obiettivi}

Il lavoro persegue sei obiettivi principali, di seguito descritti.

\paragraph{Obiettivo 1 — Caratterizzazione del drift e analisi del suo impatto.}
Il primo obiettivo consiste nell'analizzare e categorizzare le diverse
tipologie di drift che possono manifestarsi in un sistema di Machine Learning,
con particolare attenzione alla loro definizione formale e al loro impatto
sulle prestazioni dei modelli. Verranno considerate le tre principali categorie
di drift: il \textit{feature drift} (o \textit{data drift}), che riguarda il
cambiamento nella distribuzione delle variabili di ingresso; il
\textit{prediction drift}, relativo alla distribuzione delle predizioni
prodotte dal modello; e il \textit{concept drift}, che riguarda la relazione
tra input e output. La caratterizzazione includerà anche una discussione delle
modalità temporali con cui il drift può manifestarsi, ossia in modo brusco,
graduale, incrementale o ricorrente.

\paragraph{Obiettivo 2 — Metriche statistiche e metodi di detection.}
Il secondo obiettivo prevede l'analisi, la selezione e l'implementazione di
tecniche statistiche adeguate al rilevamento delle diverse forme di drift. In
questa fase vengono considerate sia tecniche classiche basate su test di
ipotesi (come il test di \textit{Kolmogorov-Smirnov}), sia tecniche pensate
specificamente per contesti di \textit{streaming} (come l'algoritmo
\textit{ADWIN}, \textit{ADaptive WINdowing}) \cite{bifet2007adwin}. Ciascun
metodo viene implementato in modo integrato all'interno di un framework
unificato, testato su dataset sintetici e validato su dataset reali provenienti
dalla letteratura di riferimento.

\paragraph{Obiettivo 3 — Progettazione e sviluppo di servizi di monitoring.}
Il terzo obiettivo consiste nella progettazione e implementazione di servizi
di monitoring capaci di tracciare il comportamento dei dati e dei modelli in
ambienti di produzione. Il servizio viene progettato secondo principi di
modularità e model-agnosticità, in modo da poter essere utilizzato con modelli
di natura differente senza richiedere modifiche al codice.

\paragraph{Obiettivo 4 — Valutazione delle strategie di retraining.}
Il quarto obiettivo prevede la sperimentazione di strategie di riaddestramento
attivate in risposta a un evento di drift o a un degrado prestazionale
osservato, con la valutazione della loro efficacia ed efficienza in scenari
controllati.

\paragraph{Obiettivo 5 — Pipeline MLOps end-to-end e containerizzazione.}
Il quinto obiettivo prevede la progettazione di una pipeline
\textit{end-to-end} che integri le componenti di monitoring e retraining
sviluppate all'interno di un framework MLOps, sfruttando strumenti come
\textit{MLflow} per il tracciamento degli esperimenti, \textit{Argo Workflows}
e \textit{Argo Events} per l'orchestrazione, e \textit{Docker} per la
containerizzazione delle componenti.

\paragraph{Obiettivo 6 — Documentazione e riproducibilità.}
L'ultimo obiettivo consiste nella redazione di una documentazione completa del
framework e delle scelte progettuali effettuate, tale da garantire la
riproducibilità degli esperimenti e la manutenibilità del codice.


\section{Contributo}
\label{sec:contributo}

Il contributo principale di questa tesi è la progettazione e l'implementazione
di un framework modulare per il drift monitoring, denominato
\texttt{drift\_framework}, che si distingue dalle soluzioni esistenti in
letteratura e sul mercato per alcuni aspetti significativi.

Un primo elemento di rilievo è la struttura architetturale, basata sul
\textit{Strategy Pattern} \cite{gof1994patterns}: le tecniche di drift
detection vengono implementate come strategie intercambiabili, che condividono
un'interfaccia comune e possono essere sostituite o affiancate senza modifiche
al resto del codice. Questa scelta consente di adottare tecniche di detection
diverse per uno stesso modello o di sperimentare tecniche nuove semplicemente
aggiungendo una nuova classe che implementi l'interfaccia astratta comune.

Un secondo elemento è la separazione netta tra \textit{strategie} — algoritmi
che operano su singoli stream numerici — e \textit{detector} — orchestratori
che operano su vettori di feature, sequenze di predizioni o stream di errori.
Questa separazione garantisce che la stessa strategia possa essere riutilizzata
per tipologie di drift diverse (feature, prediction, concept) senza
duplicazione di codice.

Un terzo elemento riguarda la \textit{model-agnosticità}: il framework non fa
alcuna assunzione sulla natura del modello di Machine Learning che monitora.
Il modello viene considerato un'entità che espone unicamente un metodo di
predizione, e tutte le informazioni necessarie al drift detection vengono
ricavate dagli stream di dati che attraversano il modello stesso. Questa scelta
rende il framework applicabile a modelli di natura fortemente eterogenea, dai
classificatori lineari alle reti neurali profonde, e a task diversi come
classificazione, regressione e clustering.

Infine, il lavoro fornisce una prima validazione sperimentale del framework
attraverso il confronto tra le due strategie implementate nella fase iniziale
del lavoro — Kolmogorov-Smirnov e ADWIN — su scenari sintetici fedeli alla
letteratura, discutendo in modo argomentato le loro proprietà complementari
ed evidenziando i \textit{trade-off} che possono guidare la scelta di una
strategia in scenari applicativi diversi.

Nel complesso, il contributo della tesi non si colloca sul piano della
reimplementazione degli algoritmi di drift detection — per i quali sono state
adottate implementazioni consolidate come \textit{scipy.stats} \cite{scipy2020}
e \textit{river} \cite{montiel2021river} — bensì sul piano della progettazione
di un livello di astrazione superiore, che permetta di comporre, sostituire ed
estendere in modo pulito i vari componenti del sistema di monitoring, aprendo
la strada a un impiego industriale sostenibile della disciplina.


\section{Struttura della tesi}
\label{sec:struttura}

La tesi è organizzata in otto capitoli.

Il presente \textbf{Capitolo~\ref{ch:introduzione}} ha introdotto il contesto
del lavoro, presentato gli obiettivi principali e riassunto il contributo
scientifico e ingegneristico che si intende offrire.

Il \textbf{Capitolo~2} è dedicato allo stato dell'arte. Verranno approfonditi
i concetti fondamentali del ciclo di vita di un modello di Machine Learning in
produzione, il fenomeno del drift nelle sue diverse tipologie, le principali
tecniche di detection presenti in letteratura e la posizione del drift
monitoring all'interno della più ampia disciplina dell'MLOps. Verrà inoltre
presentato un confronto tra le principali librerie esistenti per il drift
detection, motivando le scelte tecnologiche compiute nel lavoro.

Il \textbf{Capitolo~3} presenta l'architettura del framework sviluppato,
discutendone i principi di design, la struttura a tre livelli (strategie,
detector, servizio di monitoring) e le motivazioni delle scelte progettuali
effettuate. Particolare attenzione viene dedicata al ruolo del Strategy Pattern
come scelta architetturale abilitante e al principio Open/Closed applicato al
framework.

Il \textbf{Capitolo~4} descrive nel dettaglio le due strategie di drift
detection implementate nella fase iniziale del lavoro: il test di
Kolmogorov-Smirnov e l'algoritmo ADWIN. Per ciascuna vengono presentati i
fondamenti teorici, i dettagli implementativi, i limiti noti e le
considerazioni pratiche che ne orientano l'uso.

Il \textbf{Capitolo~5} descrive i detector e i meccanismi di orchestrazione: il
\texttt{FeatureDriftDetector} come orchestratore multi-feature attualmente
implementato, e i placeholder architetturali per il
\texttt{PredictionDriftDetector}, il \texttt{ConceptDriftDetector} e il
\texttt{DriftMonitoringService} centrale, previsti come sviluppi successivi
del lavoro.

Il \textbf{Capitolo~6} presenta la validazione sperimentale del lavoro: la
metodologia adottata, gli esperimenti condotti su dataset sintetici, i
risultati ottenuti dal confronto tra le due strategie implementate e le
successive validazioni previste su dataset reali provenienti dalla letteratura.

Il \textbf{Capitolo~7} descrive la pipeline MLOps \textit{end-to-end},
comprensiva dei meccanismi di retraining automatico e delle scelte di
integrazione con gli strumenti industriali di riferimento.

Il \textbf{Capitolo~8} chiude la tesi con un riepilogo del lavoro svolto,
un'analisi dei principali risultati ottenuti, una discussione dei limiti
identificati e la presentazione delle direzioni di sviluppo futuro.
